% arara: clean: {extensions: ['log','aux','synctex.gz','out']}
% arara: pdflatex: { shell: yes }
% arara: pdflatex: { shell: yes }
% arara: clean: {extensions: ['log','aux','synctex.gz','out']}
\documentclass[spanish,addpoints,answers,a4paper]{exam}
\usepackage[T1]{fontenc}
\usepackage[spanish]{babel}
\usepackage{libertine}
\usepackage[tmargin=2cm,rmargin=2cm,lmargin=2cm]{geometry}
\usepackage{amsmath,bm}
\usepackage[shortlabels]{enumitem}
\usepackage{siunitx}
\sisetup{
	group-four-digits = true,
	group-separator = {,}
}
\usepackage{multicol}

\usepackage{graphicx}
\graphicspath{{../../../images/}}
\pagestyle{plain}

\usepackage{minted}
\setminted[cpp]{tabsize=2}
\usepackage{hyperref}
\renewcommand{\solutiontitle}{\noindent\textbf{Solución}\par\noindent}
\renewcommand\listingscaption{Listado}
\newcommand{\unmarkedfntext}[1]{%
	\begingroup
	\renewcommand\thefootnote{}\footnote{#1}%
	\addtocounter{footnote}{-1}%
	\endgroup
}
\spanishdatedel
\decimalpoint

\def\LOGOUNI{%
	\begin{picture}(0,0)\unitlength=1cm
	\put (-13.5,-2.5) {\includegraphics[width=2.8cm]{logouni}}
	\end{picture}
}

\def\LOGOCTIC{%
	\begin{picture}(0,0)\unitlength=1cm
	\put (-8,-1.2) {\includegraphics[height=1cm]{logocticblack}}
	\end{picture}
}

\newcommand{\mychar}{%
	\begingroup\normalfont
	\includegraphics[height=\fontcharht\font`\B]{Octocat.png}%
	\endgroup
}

\begin{document}

\begin{center}
	\sffamily\bfseries\scshape
	{\Large UNIVERSIDAD NACIONAL DE INGENIERÍA}\LOGOUNI\\
	Centro de Tecnologías de la Información y Comunicaciones\LOGOCTIC\\
\end{center}

\vspace{.8cm}

\begin{center}\sffamily\bfseries\large
	Programación en \texttt{C++} \\
	TEMA $\bm{3}$: Estructuras de repetición \mintinline{cpp}{while} y \mintinline{cpp}{for}.
\end{center}

\vspace{.5cm}
\noindent
\makebox[\textwidth]{Nombres y apellidos:\enspace\hrulefill}
\vspace{0.2in}
\makebox[\textwidth]{Nombres y apellidos del instructor:\enspace MSc. Víctor Melchor Espinoza.\hfill}
%
%\section*{Ejercicios de práctica}
%
%\begin{questions}
%
%\question Solicite dos valores diferentes (no validar) e intercambiarlos.
%
%\begin{solution}
%	\begin{listing}[H]
%		\footnotesize
%		\inputminted{cpp}{1.cc}
%		\caption{Programa \texttt{1.cc}.}
%		\label{lst:2.1}
%	\end{listing}
%\end{solution}
%
%\end{questions}

\section*{Ejercicios obligatorios}

\begin{questions}

\question Elabore un programa que pida al usuario un entero positivo $n$. Luego, el programa debe pedir al usuario $n$ enteros usando un ciclo. El programa deberá encontrar el mayor, menor, y promedio de los valores dados por el usuario e imprimirlos al final.

\begin{solution}
	\begin{listing}[H]
		\footnotesize
		\inputminted{cpp}{1.cpp}
		\caption{Programa \texttt{1.cpp}.}
		\label{lst:1.9}
\end{listing}
\end{solution}

\question Desarrolle un programa que mantenga una cuenta. El programa ofrecerá al usuario un menú de opciones:

\begin{minted}{bash}
0: Limpiar la cuenta.
1: Depositar dinero a la cuenta.
2: Retirar dinero de la cuenta.
3: Mostrar el saldo actual.
4: Salir del programa.
\end{minted}

\begin{solution}
	\begin{listing}[H]
		\footnotesize
		\inputminted{cpp}{2.cpp}
		\caption{Programa \texttt{2.cpp}.}
		\label{lst:1.10}
\end{listing}
\end{solution}

\end{questions}

\begin{flushright}\bfseries
Centro de Tecnologías de la Información y Comunicaciones (CTIC)\\[2mm]
\today\unmarkedfntext{Código disponible en \href{https://github.com/carlosal1015/Cpp-Programming}{\mychar{}}.}
\end{flushright}

\end{document}